\documentclass[10pt,a4paper,landscape]{article}
\usepackage[utf8]{inputenc}
\usepackage[T1]{fontenc}
\usepackage[romanian]{babel}
\usepackage[margin=1cm]{geometry}
\usepackage{array}
\usepackage{tabularx}
\usepackage{enumitem}
\usepackage{graphicx}
\usepackage{setspace}
\usepackage{xcolor}
\usepackage{hyperref}
\usepackage{amssymb}
\usepackage{multirow}
\setlength{\parindent}{0pt}
\newcommand{\fillline}[1][6cm]{\underline{\hspace{#1}}}
\newcommand{\fillbox}[1][0.4cm]{\framebox[#1]{\rule{0pt}{#1}}}
\newcommand{\checkboxempty}{$\square$}

\begin{document}

\begin{minipage}{0.5\textwidth}
\includegraphics[height=1.3cm]{assets/stema-cluj-color.png}\hspace{0.3cm}%
\raisebox{0.4cm}{\parbox[t]{6cm}{\footnotesize
ROMÂNIA\\
Judeţul Cluj\\
Municipiul Cluj-Napoca\\
Codul de înregistrare fiscală: 4305857\\
registratura@primariaclujnapoca.ro}}
\end{minipage}\hfill
\begin{minipage}{0.45\textwidth}
\raggedleft
\footnotesize
Anexa nr. 1\\
Model 2016\\
\textbf{ITL -- 001}\\
\includegraphics[width=4cm]{assets/barcode-491007.png}
\end{minipage}

\vspace{0.3cm}

\begin{tabularx}{\textwidth}{|X|X|X|X|}
\hline
Numărul de rol nominal unic & Registrul agricol Tipul: & Volumul: & Poziţia: \\
\fillline[5cm] & \fillline[3cm] & \fillline[3cm] & \fillline[3cm] \\
\hline
\end{tabularx}

\vspace{0.3cm}

\begin{center}
\textbf{DECLARAŢIE FISCALĂ PENTRU STABILIREA IMPOZITULUI/TAXEI PE CLĂDIRILE REZIDENŢIALE/NEREZIDENŢIALE/CU DESTINAŢIE MIXTĂ AFLATE ÎN PROPRIETATEA PERSOANELOR FIZICE}
\end{center}

\textbf{I. DATE DE IDENTIFICARE A CONTRIBUABILULUI} \footnotesize (în cazul a mai mult de trei coproprietari se completează o nouă declaraţie) \normalsize

\vspace{0.2cm}

\textbf{1. Coproprietar 1}\\
\begin{tabularx}{\textwidth}{|X|X|X|X|}
\hline
Numele/Prenumele \underline{ {{nume_complet}} } & Codul de identificare fiscală & Codul numeric personal & Număr de înregistrare fiscală \\
& \fillline[4cm] & \underline{ {{cnp}} } & \fillline[4cm] \\
\hline
\end{tabularx}

\begin{tabularx}{\textwidth}{|X|X|X|X|X|}
\hline
Numărul tel./fax & Adresa de poştă electronică & Strada & Nr. & Cod poştal \\
\fillline[3cm] & \fillline[4cm] & \fillline[4cm] & \fillline[1.5cm] & \fillline[2cm] \\
\hline
Bloc & Scara \quad Etaj \quad Apartament & Localitatea & Judeţul/Sectorul & Ţara \quad Cotă proprietate \\
\fillline[1.5cm] & \fillline[3cm] & \fillline[3cm] & \fillline[3cm] & \underline{ {{cota_proprietate}} \% } \\
\hline
\end{tabularx}

\begin{tabularx}{\textwidth}{|m{2.5cm}|X|X|X|X|X|X|m{4cm}|}
\hline
\multirow{2}{2.5cm}{Adresa de corespondenţă} & Strada & Nr & Codul poştal & Blocul & Scara \quad Etaj \quad Ap. & Localitate & Sunt de acord ca actele administrative fiscale să-mi fie comunicate exclusiv la adresa de poştă electronică \\
& \fillline[2cm] & \fillline[1cm] & \fillline[1.5cm] & \fillline[1cm] & \fillline[2cm] & \fillline[2cm] & DA \checkboxempty \hspace{0.5cm} NU \checkboxempty \\
\hline
\end{tabularx}

\vspace{0.3cm}

\textbf{2. Coproprietar 2}\\
\begin{tabularx}{\textwidth}{|X|X|X|X|}
\hline
Numele \fillline[5cm] & Codul de identificare fiscală & Codul numeric personal & Număr de înregistrare fiscală \\
Prenumele \fillline[5cm] & \fillline[4cm] & \fillline[4cm] & \fillline[4cm] \\
\hline
\end{tabularx}

\begin{tabularx}{\textwidth}{|X|X|X|X|X|}
\hline
Numărul tel./fax & Adresa de poştă electronică & Strada & Nr. & Cod poştal \\
\fillline[3cm] & \fillline[4cm] & \fillline[4cm] & \fillline[1.5cm] & \fillline[2cm] \\
\hline
Bloc & Scara \quad Etaj \quad Apartament & Localitatea & Judeţul/Sectorul & Ţara \quad Cotă proprietate \\
\fillline[1.5cm] & \fillline[3cm] & \fillline[3cm] & \fillline[3cm] & \fillline[3cm] \\
\hline
\end{tabularx}

\begin{tabularx}{\textwidth}{|m{2.5cm}|X|X|X|X|X|X|m{4cm}|}
\hline
\multirow{2}{2.5cm}{Adresa de corespondenţă} & Strada & Nr & Codul poştal & Blocul & Scara \quad Etaj \quad Ap. & Localitate & Sunt de acord ca actele administrative fiscale să-mi fie comunicate exclusiv la adresa de poştă electronică \\
& \fillline[2cm] & \fillline[1cm] & \fillline[1.5cm] & \fillline[1cm] & \fillline[2cm] & \fillline[2cm] & DA \checkboxempty \hspace{0.5cm} NU \checkboxempty \\
\hline
\end{tabularx}

\vspace{0.3cm}

\textbf{3. Coproprietar 3}\\
\begin{tabularx}{\textwidth}{|X|X|X|X|}
\hline
Numele \fillline[5cm] & Codul de identificare fiscală & Codul numeric personal & Număr de înregistrare fiscală \\
Prenumele \fillline[5cm] & \fillline[4cm] & \fillline[4cm] & \fillline[4cm] \\
\hline
\end{tabularx}

\begin{tabularx}{\textwidth}{|X|X|X|X|X|}
\hline
Numărul tel./fax & Adresa de poştă electronică & Strada & Nr. & Cod poştal \\
\fillline[3cm] & \fillline[4cm] & \fillline[4cm] & \fillline[1.5cm] & \fillline[2cm] \\
\hline
Bloc & Scara \quad Etaj \quad Apartament & Localitatea & Judeţul/Sectorul & Ţara \quad Cotă proprietate \\
\fillline[1.5cm] & \fillline[3cm] & \fillline[3cm] & \fillline[3cm] & \fillline[3cm] \\
\hline
\end{tabularx}

\begin{tabularx}{\textwidth}{|m{2.5cm}|X|X|X|X|X|X|m{4cm}|}
\hline
\multirow{2}{2.5cm}{Adresa de corespondenţă} & Strada & Nr & Codul poştal & Blocul & Scara \quad Etaj \quad Ap. & Localitate & Sunt de acord ca actele administrative fiscale să-mi fie comunicate exclusiv la adresa de poştă electronică \\
& \fillline[2cm] & \fillline[1cm] & \fillline[1.5cm] & \fillline[1cm] & \fillline[2cm] & \fillline[2cm] & DA \checkboxempty \hspace{0.5cm} NU \checkboxempty \\
\hline
\end{tabularx}

\vspace{0.3cm}

\textbf{II. DATE DE IDENTIFICARE A ÎMPUTERNICITULUI} \footnotesize (Împuternicirea nu este transmisibilă şi încetează la data revocării de către contribuabil sau la data decesului acestuia.) \normalsize

\begin{tabularx}{\textwidth}{|X|X|}
\hline
Numele \fillline[6cm] & Codul numeric personal \fillline[4cm] \\
\hline
Prenumele \fillline[6cm] & \\
\hline
\end{tabularx}

\begin{tabularx}{\textwidth}{|X|X|X|}
\hline
Numărul tel./fax \fillline[3cm] & Adresa de poştă electronică \fillline[4cm] & Strada \fillline[5cm] \\
\hline
\end{tabularx}

\begin{tabularx}{\textwidth}{|X|X|X|X|X|X|X|}
\hline
Număr & Cod poştal & Blocul & Scara \quad Etaj \quad Apartament & Localitate & Judeţul/Sectorul & Ţara \\
\fillline[1.5cm] & \fillline[2cm] & \fillline[1.5cm] & \fillline[3cm] & \fillline[3cm] & \fillline[3cm] & \fillline[2cm] \\
\hline
\end{tabularx}

\vspace{0.3cm}

\textbf{III. Adresa de reziden\c{t}ă a contribuabilului în alt stat decât România}\\
\fillline[18cm]

\vspace{0.3cm}

\textbf{IV. DATELE CLĂDIRII NECESARE STABILIRII IMPOZITULUI/TAXEI DATORAT/E PE CLĂDIRI REZIDENŢIALE}

\textbf{Tip operaţiune:} \underline{ {{tip_operatiune}} } \hspace{1cm} \textbf{Adresă imobil:} \underline{ {{adresa_imobil}} } \hspace{1cm} \textbf{Nr. cadastral:} \underline{ {{nr_cadastral}} }

\vspace{0.2cm}

\begin{tabularx}{\textwidth}{|X|X|X|X|X|}
\hline
Nr. act dobândire: & Data dobândirii: & Valoare achiziţie (lei) & Cota: & În calitate de: \\
\fillline[3cm] & \underline{ {{data_dobandirii}} } & \fillline[3cm] & \underline{ {{cota_proprietate}} \% } & \fillline[3cm] \\
\hline
\end{tabularx}

Beneficiez de scutire/reducere de la plata impozitului pe clădire \quad DA \checkboxempty \quad NU \checkboxempty

\begin{tabularx}{\textwidth}{|X|X|X|X|X|X|X|X|}
\hline
Strada & Nr & Bloc & Scara & Etaj & Ap & An construire & bloc cu mai mult de 3 etaje şi 8 apartamente DA \checkboxempty \quad NU \checkboxempty \\
\fillline[3cm] & \fillline[1cm] & \fillline[1cm] & \fillline[1cm] & \fillline[1cm] & \fillline[1cm] & \fillline[2cm] & \\
\hline
\end{tabularx}

\vspace{0.3cm}

\textbf{Tabel rezidenţial -- tipuri de clădiri / construcţii anexe / subsol-demisol-mansardă}

\renewcommand{\arraystretch}{1.3}
\begin{tabularx}{\textwidth}{|m{2.6cm}|X|X|X|X|X|X|}
\hline
\textbf{Specificaţii} & \textbf{A.} Clădire cu cadre din beton armat sau cu pereţi exteriori din cărămidă arsă sau din orice alte materiale rezultate în urma unui tratament termic şi/sau chimic & \textbf{B.} Clădire cu pereţii exteriori din lemn, din piatră naturală, din cărămidă nearsă, din vălătuci sau din orice alte materiale nesupuse unui tratament termic şi/sau chimic & \textbf{C.} Clădire-anexă cu cadre din beton armat sau cu pereţi exteriori din cărămidă arsă sau din orice alte materiale rezultate în urma unui tratament termic şi/sau chimic & \textbf{D.} Clădire-anexă cu pereţii exteriori din lemn, din piatră naturală, din cărămidă nearsă, din vălătuci sau din orice alte materiale nesupuse unui tratament termic şi/sau chimic & \textbf{E.} În cazul contribuabilului care deţine la aceeaşi adresă încăperi amplasate la subsol, demisol şi/sau la mansardă, utilizate ca locuinţă, în oricare dintre tipurile de clădiri prevăzute la lit. A-D & \textbf{F.} În cazul contribuabilului care deţine la aceeaşi adresă încăperi amplasate la subsol, la demisol şi/sau la mansardă, utilizate în alte scopuri decât cel de locuinţă, în oricare dintre tipurile de clădiri prevăzute la lit. A-D \\
\hline
Condiţii de dotare cu instalaţii de apă, canalizare, electrice, încălzire (condiţii cumulative) & instalaţii de apă \checkboxempty \quad Da \checkboxempty \quad Nu \checkboxempty \newline canalizare \checkboxempty \quad Da \checkboxempty \quad Nu \checkboxempty \newline electrice \checkboxempty \quad Da \checkboxempty \quad Nu \checkboxempty \newline încălzire \checkboxempty \quad Da \checkboxempty \quad Nu \checkboxempty & instalaţii de apă \checkboxempty \quad Da \checkboxempty \quad Nu \checkboxempty \newline canalizare \checkboxempty \quad Da \checkboxempty \quad Nu \checkboxempty \newline electrice \checkboxempty \quad Da \checkboxempty \quad Nu \checkboxempty \newline încălzire \checkboxempty \quad Da \checkboxempty \quad Nu \checkboxempty & instalaţii de apă \checkboxempty \quad Da \checkboxempty \quad Nu \checkboxempty \newline canalizare \checkboxempty \quad Da \checkboxempty \quad Nu \checkboxempty \newline electrice \checkboxempty \quad Da \checkboxempty \quad Nu \checkboxempty \newline încălzire \checkboxempty \quad Da \checkboxempty \quad Nu \checkboxempty & instalaţii de apă \checkboxempty \quad Da \checkboxempty \quad Nu \checkboxempty \newline canalizare \checkboxempty \quad Da \checkboxempty \quad Nu \checkboxempty \newline electrice \checkboxempty \quad Da \checkboxempty \quad Nu \checkboxempty \newline încălzire \checkboxempty \quad Da \checkboxempty \quad Nu \checkboxempty & instalaţii de apă \checkboxempty \quad Da \checkboxempty \quad Nu \checkboxempty \newline canalizare \checkboxempty \quad Da \checkboxempty \quad Nu \checkboxempty \newline electrice \checkboxempty \quad Da \checkboxempty \quad Nu \checkboxempty \newline încălzire \checkboxempty \quad Da \checkboxempty \quad Nu \checkboxempty & instalaţii de apă \checkboxempty \quad Da \checkboxempty \quad Nu \checkboxempty \newline canalizare \checkboxempty \quad Da \checkboxempty \quad Nu \checkboxempty \newline electrice \checkboxempty \quad Da \checkboxempty \quad Nu \checkboxempty \newline încălzire \checkboxempty \quad Da \checkboxempty \quad Nu \checkboxempty \\
\hline
Suprafaţa utilă (m\textsuperscript{2}) & {{suprafata}} & & & & & \\
\hline
Suprafaţa construită desfăşurată (m\textsuperscript{2}) & & & & & & \\
\hline
\end{tabularx}
\renewcommand{\arraystretch}{1}

\footnotesize\textit{Tabelul de mai sus se completează inclusiv pentru clădirile cu destinaţie agricolă}\normalsize

\vspace{0.3cm}

\textbf{V. DATELE CLĂDIRII NECESARE STABILIRII IMPOZITULUI/TAXEI DATORAT/E PE CLĂDIRI NEREZIDENŢIALE aflate în proprietatea persoanelor fizice}

\begin{tabularx}{\textwidth}{|X|X|X|X|X|X|X|X|}
\hline
Strada & Nr & Bloc & Scara & Etaj & Ap & An construire & bloc cu mai mult de 3 etaje şi 8 apartamente DA \checkboxempty \quad NU \checkboxempty \\
\fillline[3cm] & \fillline[1cm] & \fillline[1cm] & \fillline[1cm] & \fillline[1cm] & \fillline[1cm] & \fillline[2cm] & \\
\hline
\end{tabularx}

Suprafaţa clădirii: utilă \fillline[3cm] m\textsuperscript{2} \hspace{1cm} construită desfăşurată \fillline[3cm] m\textsuperscript{2}

\vspace{0.3cm}

\begin{tabularx}{\textwidth}{|X|X|X|}
\hline
\textbf{a)} raport de evaluare întocmit de un evaluator autorizat în ultimii 5 ani anteriori anului de referinţă & \textbf{b)} clădiri nou construite, în ultimii 5 ani anteriori anului de referinţă & \textbf{c)} clădiri dobândite în ultimii 5 ani anteriori anului de referinţă \\
\hline
Valoarea clădirii (lei) \fillline[4cm] & Valoarea finală a lucrărilor de construcţii (lei) \fillline[3cm] & Valoarea din actul de proprietate (lei) \fillline[3cm] \\
\hline
Dată întocmire raport \fillline[3cm] & Dată întocmire proces verbal de recepţie \fillline[3cm] & Dată încheiere act de dobândire \fillline[3cm] \\
\hline
\end{tabularx}

\vspace{0.2cm}

\footnotesize\textit{Notă: în cazul în care contribuabilul persoană fizică nu deţine informaţiile necesare completării punctului V completează informaţiile de la punctul IV şi consemnează în cele ce urmează olograf faptul că utilizează clădirea în scop nerezidenţial:}\normalsize

\vspace{0.3cm}
\hrulefill \\
\hrulefill

\vspace{0.5cm}

\textbf{VI.} În cazul în care clădirea este utilizată în scop mixt (rezidenţial şi nerezidenţial), iar \textbf{SUPRAFEŢELE POT FI DELIMITATE} atunci se completează în mod corespunzător pct. IV, respectiv V de mai sus.

\vspace{0.3cm}

\textbf{VII.} În cazul în care clădirea este utilizată în scop mixt (rezidenţial şi nerezidenţial), iar \textbf{SUPRAFEŢELE NU POT FI DELIMITATE} atunci contribuabilul bifează situaţia în care se află:

\vspace{0.2cm}

\checkboxempty\ \textbf{a)} La adresa poştală a clădirii este înregistrat un domiciliu fiscal al unui operator economic, dar nu se desfăşoară nici o activitate economică.

\textit{În acest caz, contribuabilul completează punctul IV şi consemnează olograf pe proprie răspundere acest fapt pe rândul următor:}\\
\hrulefill

\vspace{0.2cm}

\checkboxempty\ \textbf{b)} La adresa poştală a clădirii este înregistrat un domiciliu fiscal al unui operator economic care desfăşoară activitate economică, iar cheltuielile cu utilităţile sunt în sarcina operatorului economic. În acest caz, contribuabilul completează punctul V.

\vspace{0.5cm}

Anexez la prezenta declaraţie documentele, certificate de conformitate cu originalul, conform art. 64 alin. (5) din Legea nr. 207/2015 privind Codul de procedură fiscală, cu modificările şi completările ulterioare, după cum urmează:\textsuperscript{1}

\begin{tabularx}{\textwidth}{XXX}
1. \fillline[5cm] & 4. \fillline[5cm] & 7. \fillline[5cm] \\
2. \fillline[5cm] & 5. \fillline[5cm] & 8. \fillline[5cm] \\
3. \fillline[5cm] & 6. \fillline[5cm] & 9. \fillline[5cm] \\
\end{tabularx}

\vspace{0.3cm}

\textbf{Sub sancţiunile aplicate faptei de fals în acte publice, declar că:}

1. datele înscrise în prezentul formular, precum şi orice documente anexate depuse de mine sunt corecte şi complete, conforme cu realitatea;

2. în cazul în care intervin modificări privind situaţia juridică a contribuabilului ori a bunului impozabil/taxabil, mă oblig să depun o nouă declaraţie fiscală care să reflecte realitatea, în termen de 30 de zile de la apariţia situaţiei respective.

\vspace{0.3cm}

\begin{tabularx}{\textwidth}{|X|X|X|X|X|}
\hline
Posed actul de identitate & TIPUL \fillline[2cm] & seria \fillline[2cm] & numărul \fillline[3cm] & eliberat de \fillline[3cm] \quad La data de \fillline[2.5cm] \\
\hline
\end{tabularx}

\vspace{0.5cm}

\begin{tabularx}{\textwidth}{|X|X|}
\hline
Semnătura contribuabilului & Semnătura împuternicitului \\
& \\
\underline{ {{nume_complet}} } & \fillline[7cm] \\
Data semnării \underline{ {{data_dobandirii}} } & Data semnării \fillline[4cm] \\
\hline
\end{tabularx}

\vspace{0.5cm}

\footnotesize\textit{În cazul în care contribuabilul nu poate citi sau scrie, declaraţia fiscală se completează de o persoană agreată de către acesta, care îi va citi integral conţinutul declaraţiei fiscale şi va semna pentru conformitate.}\normalsize

\vspace{0.3cm}

Date de identificare ale acestei persoane:

\begin{tabularx}{\textwidth}{|X|X|X|X|}
\hline
Posed actul de identitate TIPUL \fillline[2cm] & seria \fillline[2cm] & numărul \fillline[3cm] & eliberat de \fillline[3cm] \\
\hline
\end{tabularx}

\vspace{0.3cm}

Data şi semnătura \fillline[8cm]

\vspace{0.5cm}

\textit{Prin semnarea prezentei am luat la cunostinta ca declararea necorespunzatoare a adevarului se pedepseste conform legii penale, cele declarate fiind corecte si complete.}

\vspace{0.5cm}

\checkboxempty\ Prin prezenta solicit comunicarea răspunsului pe următoarea adresă de e-mail: \fillline[10cm]

\checkboxempty\ Mă oblig să comunic instituţiei orice modificare intervine în legătură cu această adresă de e-mail

\checkboxempty\ Îmi exprim consimţământul ca Primăria Municipiului Cluj-Napoca să comunice orice informaţii, date personale, clarificări şi completări pe adresa de e-mail indicată mai sus

\checkboxempty\ Am luat la cunoştinţă faptul că în cazul nefuncţionării serverului de e-mail comunicat sau în cazul adresei greşite de e-mail, Municipiul Cluj-Napoca nu poate fi tras la răspundere pentru acest lucru

\vspace{0.5cm}

\textbf{În acest tip de chenar completează organul fiscal local:}

\begin{tabularx}{\textwidth}{|X|X|}
\hline
Prenumele şi numele \fillline[6cm] & Rangul localităţii \fillline[4cm] \\
\hline
Semnătura \fillline[6cm] & Zona în cadrul localităţii \fillline[4cm] \\
\hline
Data semnării \fillline[4cm] & \\
\hline
\end{tabularx}

\vspace{0.5cm}

\footnotesize
\textsuperscript{1}Primăria municipiului Cluj-Napoca asigură gratuit fotocopierea documentelor.

Pentru documentele emise de către alte instituţii publice, organe de specialitate ale administraţiei publice centrale şi locale, precum şi persoane juridice de drept privat, care potrivit legii au obţinut statut de utilitate publică sau sunt autorizate să presteze un serviciu public, în regim de putere publică, pe care solicitantul nu le depune, acesta are posibilitatea de a-şi da consimţământul expres ca Primăria municipiului Cluj-Napoca să obţină în numele său copii sau extrase ale acestora.

\vspace{0.3cm}

Timp estimativ de completare: 10 minute

Datele personale care vă sunt solicitate prin prezenta cerere vor fi prelucrate numai în vederea procesării şi soluţionării solicitării dumneavoastră. Primăria municipiului Cluj-Napoca garantează securitatea procesării datelor şi arhivarea acestora în conformitate cu prevederile legale în vigoare. Responsabilul Primăriei municipiului Cluj-Napoca cu protecţia datelor poate fi contactat pe adresa de email dpo@primariaclujnapoca.ro.

În conformitate cu Regulamentul nr. 679/2016 aveţi dreptul de a solicita Primăriei municipiului Cluj-Napoca, în ceea ce priveşte datele cu caracter personal referitoare la persoana vizată, accesul la acestea, rectificarea sau ştergerea acestora sau restricţionarea prelucrării sau a dreptului de a vă opune prelucrării, precum şi a dreptului la portabilitatea datelor.

\end{document}
